\documentclass[10pt,a4paper]{amsart}
\usepackage[margin=19mm]{geometry}
\usepackage{amsmath,amssymb,mathtools}
\usepackage[scr=rsfs,cal=euler]{mathalfa}
\usepackage{xfrac}
\usepackage[unicode,hidelinks,bookmarks=false]{hyperref}
\newcommand{\ie}{i.e.\ }
\newcommand{\cf}{cf.\ }
\newcommand{\resp}{respectively\ }
\newcommand{\Subsection}{\subsection}
\providecommand{\nobreakdash}{\nobreak\hbox{-}}
\setlength{\emergencystretch}{2em}
\multlinegap0pt
\usepackage[shortlabels]{enumitem}
\usepackage{amssymb}
\usepackage{mathtools}
\usepackage{xcolor}
\usepackage{algorithm}\usepackage{algpseudocode}
\newtheorem{assumption}{Assumption}
\newtheorem{theorem}{Theorem}
\newtheorem{lemma}{Lemma}
\newcommand{\E}{\mathbb{E}}
\renewcommand{\P}{\mathbb{P}}
\newcommand{\R}{\mathbb{R}}
\newcommand{\eps}{\epsilon}
\newcommand{\norm}[1]{\left\|{#1}\right\|} % A norm with 1 argument
\newcommand{\profile}{\psi}
\newcommand{\supp}{\mathop{\rm supp}} 
\title{A Compositional Kernel Model for Feature Learning}
\author{Feng Ruan, Keli Liu, Michael Jordan}
\date{Source: arXiv:2509.14158v3 (CC BY 4.0)}
\begin{document}
\maketitle

\noindent\emph{Source and licence.} A Compositional Kernel Model for Feature Learning. Version arXiv:2509.14158v3 (CC BY 4.0), distributed by arXiv under the Creative Commons Attribution 4.0 International licence, http://creativecommons.org/licenses/by/4.0/, as recorded in the arXiv licence metadata for this exact version and shown on the abstract page https://arxiv.org/abs/2509.14158v3. Full licence notice: \texttt{licenses/arxiv-2509.14158-CC-BY-4.0.txt}.\par\smallskip

\noindent\textit{Editorial scope.} Counted unit: the article's theorem theorem:uniform-lower-bound-on-M-i-beta (source0.tex lines 2529--2539). The article's complete original proof of it is delivered with it (source0.tex lines 2548--2643, one \texttt{proof} environment of 96 lines, which the article places away from the statement). No mathematics is written, repaired or re-derived: the statement and the proof are the article's own lines, in the article's own notation, and no retained line is substituted or re-flowed.  Retained uncounted as the source-local context the proof needs: lines 1138--1141; lines 1155--1159; lines 2356--2359; lines 2419--2431.  Nothing was omitted from any retained span, so the package carries no omission record.  No retained line contains a citation, so the wrapper carries no bibliography and no bibliography entry.  No retained line contains a figure environment, an image-inclusion command or drawing code, so no image is delivered and the package declares no asset dependency.  Editorial formatting only, in addition: an AMS \texttt{amsart} preamble that restates the article's own declarations and macros used by the retained lines, namely the environments \texttt{assumption}, \texttt{theorem}, \texttt{lemma} on separate counters, together with the standard packages those lines need.  The retained environments print in this document as Assumption 1, Theorem 1, Lemma 1 in order of appearance, whereas the article prints the same items with the numbering of its own full document.  The complete native source of the article as distributed in the arXiv e-print for this exact version is bundled unchanged as \texttt{sources/185218/source0.tex}.
\begin{assumption}
\label{assumption:existence-of-csfs}
The minimal sufficient feature set $S_*$ exists for the distribution $(X, Y) \sim \P$. 
\end{assumption} 

\begin{assumption}
\label{assumption:independence}
Assume that $X_{S_*}$ and $X_{S_*^c}$ are probabilistically independent. Here, $S_*$ denotes the 
minimal sufficient feature set of the distribution $(X, Y) \sim \P$, and $S_*^c = \{1, 2, \ldots, d\} \backslash S_*$ denotes its complement.
\end{assumption} 

\begin{assumption}
\label{assumption:feature-variable}
The variables $(X_j)_{j \in S_*}$ in $S_*$ are mutually independent.
\end{assumption} 

\begin{theorem} 
\label{theorem:gradient-expression-ANOVA}
Let $\beta \in \R^d$ satisfy $\beta_i = 0$, and let $\lambda \in (0, 1)$. Assume Assumptions~\ref{assumption:existence-of-csfs},~\ref{assumption:independence},~\ref{assumption:feature-variable}. Let 
\begin{equation*} 
	 \mathcal{M}_i(\beta): =  \iiint \left|\mathbb{E}\left[(Y-\E[Y|X_{\supp(\beta)}]) e^{-2\pi i \langle \omega, \beta \circ X \rangle}  e^{-2\pi i \zeta X_i}\right]\right|^2 \frac{1}{2\pi^2 \zeta^2} d\zeta \cdot t q_t(\omega) \, d\omega  \mu(dt).
\end{equation*}
Then
\begin{equation*}
	\mathrm{D}\mathcal{J}(\beta, \lambda)[e_i] = - \frac{1}{\lambda} \mathcal{M}_i(\beta) + \frac{1}{\sqrt{\lambda}}\mathcal{R}_i(\beta, \lambda)
\end{equation*}
where $|\mathcal{R}_i(\beta, \lambda)| \le C$ 
for some constant $C < \infty$ depending only on $M_X$, $M_Y$, and $M_\mu$.
\end{theorem} 

\begin{theorem}
\label{theorem:uniform-lower-bound-on-M-i-beta}
Assume Assumptions~\ref{assumption:existence-of-csfs},~\ref{assumption:independence},~\ref{assumption:feature-variable}. Suppose $i \in S_*$ has a main effect: $\E[Y|X_i] \ne 0$. Then, for every compact set $\mathsf{C} \subset \{\beta \in \R^d : \beta_i = 0\}$, we have
	\begin{equation*}
		\inf_{\beta \in \mathsf{C}} \mathcal{M}_i(\beta) > 0.
	\end{equation*}
Consequently, there exist constants $\varepsilon > 0$ and $\lambda_0 > 0$ such that for all $\lambda \in (0, \lambda_0)$ and all $\beta \in \mathsf{C}$,
\begin{equation*}
	\lambda \cdot \mathrm{D}\mathcal{J}(\beta, \lambda)[e_i] \le -\varepsilon.
\end{equation*}
\end{theorem}

\begin{proof}[Proof of Theorem~\ref{theorem:uniform-lower-bound-on-M-i-beta}]
We note $\mathcal{M}_i(\beta)$ is not continuous in $\beta$ due to the discontinuity of $\supp(\beta)$. 
To remedy this issue, we define for every $i$, subset $A \subset \{1, 2, \ldots, d\}$, and $\beta \in \R^d$
\begin{equation*}
	 \mathcal{N}_{i, A}(\beta) = \iiint \left|\E[(Y-\E[Y|X_A]) e^{-2\pi i \langle \omega, \beta\circ X \rangle} e^{-2\pi i \zeta X_i}] \right|^2 \frac{1}{2\pi^2 \zeta^2} d\zeta \cdot t q_t(\omega) \, d\omega  \mu(dt).
\end{equation*} 
By definition, $\mathcal{M}_i(\beta)  = \mathcal{N}_{i, A}(\beta)$ when $\supp(\beta) = A$. 

We claim that $\mathcal{N}_{i, A}(\beta)$ is lower semicontinuous. Fix any sequence $\beta^{(n)} \to \beta$. For each fixed $(\omega, \zeta, t)$, the inner expectation in the integrand is continuous in $\beta$, so the integrand converges pointwise. Moreover, the integrand is nonnegative. By Fatou’s lemma, we conclude
\[
\liminf_{n \to \infty} \mathcal{N}_{i, A}(\beta^{(n)}) \ge \mathcal{N}_{i, A}(\beta),
\]
so $\mathcal{N}_{i, A}$ is lower semicontinuous on $\R^d$. 

Next, we show $\mathcal{N}_{i, A}(\beta)  > 0$ when $i \not \in A$. To see this, 
suppose $ \mathcal{N}_{i, A}(\beta) = 0$. Then the integrand vanishes almost everywhere, 
\begin{equation*}
	\mathbb{E}[(Y-\E[Y|X_A]) e^{-2\pi i \langle \omega, \beta \circ X \rangle}  e^{-2\pi i \zeta X_i}] = 0~~a.e.~\omega, \zeta.
\end{equation*}
Let $B:=\supp(\beta)$. Then
\begin{equation*}
	\E\left[\E[(Y-\E[Y|X_A])|X_{B \cup \{i\}}]e^{-2\pi i \langle \omega, \beta \circ X \rangle}  e^{-2\pi i \zeta X_i}\right] = 0~a.e.~\omega, \zeta
\end{equation*} 
which implies the random variable $\E[(Y-\E[Y|X_A])|X_{B \cup \{i\}}]$
must be zero almost surely. However, consider its
inner product with the main effect component $f_{{i}}(X_i)$. If $i \not\in A$, then 
\begin{equation*}
\begin{split}
\mathbb{E}\left[\mathbb{E}[Y - \mathbb{E}[Y | X_A] | X_{B \cup \{i\}}] \cdot f_{{i}}(X_i)\right]
&= \mathbb{E}[(Y - \mathbb{E}[Y | X_A]) \cdot f_{{i}}(X_i)] \\
&= \mathbb{E}[Y \cdot f_{{i}}(X_i)] - \mathbb{E}[\mathbb{E}[Y| X_A] \cdot f_{{i}}(X_i)] \\
&= \mathbb{E}[f_{{i}}^2(X_i)],
\end{split}
\end{equation*}
where the last equality uses that $\E[f_i(X_i)] = 0$ and $X_i$ is independent of $X_A$ whenever $i \notin A$. Since $f_i(X_i) = \E[Y | X_i] \not\equiv 0$ by assumption, we have $\E[f_i^2(X_i)] > 0$, contradicting the claim that the random variable is zero. Thus, 
we have $\mathcal{N}_{i, A}(\beta)  > 0$ when $i \not \in A$.


%\vspace{.5em} 
%\begin{lemma}
%If $\supp(\beta) \subseteq A$, then 
%\begin{equation*}
%	 \mathcal{N}_{i, A}(\beta) = 
%	 	\E[(Y- \E[Y|X_{A}])(Y'- \E[Y'|X'_{A}])
%		\profile^\prime(\norm{\beta \circ (X-X')}_{\ell_1}) |X_i - X_i'|].
%\end{equation*}
%In particular, $\beta \mapsto  \mathcal{N}_{i, A}(\beta)$ is continuous 
%on the closed set $\{\beta \in \R^d: \supp(\beta) \subseteq A\}$.
%\end{lemma} 
%\begin{proof}
%	At $\beta$ where $\supp(\beta) \subseteq A$, 
%		\begin{equation*}
%			\E[(Y-\E[Y|X_A])|X_{\supp(\beta)}] = 0.
%		\end{equation*} 
%	The identity then follows from Lemma xxx. For continuity, note that $\profile^\prime$ is continuous and 
%	bounded by $|\profile^\prime_+(0)|$, so the integrand is dominated by $|\profile^\prime_+(0)| |YY'| |X_i - X_i'|$ 
%	which is integrable. 
%	By the dominated convergence theorem, $\mathcal{N}_{i, A}(\beta)$ is continuous on $\{\beta: \supp(\beta) \subseteq A\}$.
%\end{proof} 

We now return to the proof of Theorem~\ref{theorem:uniform-lower-bound-on-M-i-beta}. Fix any subset 
$A \subset \{1, \dots, d\} \setminus \{i\}$.  The function $\mathcal{N}_{i, A}(\beta)$ is lower semicontinuous 
and  strictly positive on the compact set $\mathsf{C}$. By the extreme value theorem, it attains its minimum
on $\mathsf{C}$. Thus, there exists a constant $\varepsilon_A > 0$ such that
\[
\mathcal{N}_{i, A}(\beta) \ge \varepsilon_A \quad \forall \beta \in \mathsf{C}.
\]
Let $\varepsilon > 0$ be the minimum of all such $\varepsilon_A$ over subsets $A \subset \{1, \dots, d\} \setminus \{i\}$. 

For any $\beta \in \mathsf{C}$, we have $i \notin \supp(\beta)$ by assumption, so
\begin{equation*}
	\mathcal{M}_i(\beta) = \mathcal{N}_{i, \supp(\beta)}(\beta) \ge \eps_{\supp(\beta)} \ge \eps > 0.
\end{equation*}
This shows $\inf_{\beta \in \mathsf{C}} \mathcal{M}_i(\beta)  > 0$, as desired.
%\end{proof} 
%
%
%\begin{theorem}
%Assume Assumptions~\ref{assumption:existence-of-csfs},~\ref{assumption:independence},~\ref{assumption:feature-variable}. Suppose $i \in S_*$ has main effect: $\E[Y|X_i] \ne 0$.  
%For any compact set $\mathsf{C} \subset \{\beta \in \R^d : \beta_i = 0\}$, there exists $\eps > 0$ and $\lambda_0$ such that 
%for all $\lambda \in (0, \lambda_0)$
% \begin{equation*}
% 	\lambda \cdot \mathrm{D}\mathcal{J}(\beta, \lambda)[e_i] \le -\eps,~~~\forall \beta \in \mathsf{C}.
% \end{equation*}
%\end{theorem}
%
%\begin{proof} 

Finally, by Theorem~\ref{theorem:gradient-expression-ANOVA}, there exists $C < \infty$ such that for all $\lambda \in (0, 1)$,
and for all $\beta \in \mathsf{C}$
\begin{equation*}
	\lambda \cdot \mathrm{D}\mathcal{J}(\beta, \lambda)[e_i] \le -\mathcal{M}_i(\beta) + C \sqrt{\lambda}.
\end{equation*}
It follows that for sufficiently small \( \lambda > 0 \), we have
$\lambda \cdot \mathrm{D}\mathcal{J}(\beta, \lambda)[e_i] \le -\eps$ must hold for all $\beta \in \mathsf{C}$.
\end{proof} 

\end{document}
